\chapter{Serial COM1}
\label{cap:serial}

\section{Vis\~ao geral}

O m\'odulo \code{src/serial.rs} implementa um subconjunto do driver 8250 do
Linux para a UART 16550 em COM1. A serial tem dois pap\'eis:

\begin{itemize}
  \item \textbf{depura\c{c}\~ao}: todo texto impresso no VGA tamb\'em \'e
        enviado para COM1;
  \item \textbf{entrada}: al\'em do teclado PS/2, o shell aceita comandos
        digitados no terminal serial, o que viabiliza testes automatizados.
\end{itemize}

N\~ao h\'a interrup\c{c}\~oes, FIFO baseada em IRQ, \emph{termios}, \emph{tty}
layer ou controle de fluxo. O acesso \'e feito por polling de registradores
de I/O.

\section{Hardware}

A UART 16550 em COM1 usa uma base de 8 portas a partir de \code{0x3F8}:

\begin{table}[H]
\centering
\caption{Registradores da UART 16550 (COM1).}
\label{tab:uartregs}
\begin{tabular}{@{}llll@{}}
\toprule
\textbf{Offset} & \textbf{Nome} & \textbf{Leitura} & \textbf{Escrita} \\
\midrule
\code{+0} & RBR / THR & Receive Buffer & Transmit Holding \\
\code{+1} & IER / DLM & Interrupt Enable & Divisor Latch High \\
\code{+2} & IIR / FCR & Interrupt ID & FIFO Control \\
\code{+3} & LCR & Line Control & Line Control \\
\code{+4} & MCR & Modem Control & Modem Control \\
\code{+5} & LSR & Line Status & --- \\
\code{+6} & MSR & Modem Status & --- \\
\code{+7} & SCR & Scratch & Scratch \\
\bottomrule
\end{tabular}
\end{table}

O registrador LSR tem os bits usados pelo driver:

\begin{table}[H]
\centering
\caption{Bits do Line Status Register (LSR).}
\label{tab:lsr}
\begin{tabular}{@{}lll@{}}
\toprule
\textbf{Bit} & \textbf{M\'ascara} & \textbf{Significado} \\
\midrule
0 & \code{0x01} & Dado dispon\'ivel para leitura \\
5 & \code{0x20} & Transmitter Holding Register vazio (pode enviar) \\
\bottomrule
\end{tabular}
\end{table}

\section{Inicializa\c{c}\~ao}

A configura\c{c}\~ao implementada \'e 38400 baud, 8 bits de dados, sem
paridade e 1 bit de parada (8N1), com FIFO habilitada e IRQs desligadas.

\begin{table}[H]
\centering
\caption{Sequ\^encia de inicializa\c{c}\~ao da COM1.}
\label{tab:uartinit}
\begin{tabular}{@{}llp{7.2cm}@{}}
\toprule
\textbf{Offset} & \textbf{Valor} & \textbf{Efeito} \\
\midrule
\code{+1} & \code{0x00} & Desabilita interrup\c{c}\~oes. \\
\code{+3} & \code{0x80} & Liga DLAB para acessar o divisor. \\
\code{+0} & \code{0x03} & Divisor baixo = 3. \\
\code{+1} & \code{0x00} & Divisor alto = 0. \\
\code{+3} & \code{0x03} & 8 bits, sem paridade, 1 stop; desliga DLAB. \\
\code{+2} & \code{0xC7} & FIFO habilitada, limpa buffers, gatilho de 14 bytes. \\
\code{+4} & \code{0x0B} & DTR, RTS e OUT2; IRQs desligadas. \\
\bottomrule
\end{tabular}
\end{table}

O divisor 3 produz:

\[
\text{baud} = \frac{115200}{3} = 38400.
\]

Ap\'os a configura\c{c}\~ao, o driver l\^e o LSR. Se o valor for \code{0xFF},
considera que n\~ao h\'a UART no barramento e marca \code{PRESENT = false}.
Essa heur\'istica funciona porque um barramento sem dispositivo geralmente
devolve todos os bits em 1.

\begin{nota}[Por que 38400 e n\~ao 115200?]
O divisor 3 \'e um valor conservador e amplamente suportado por emula\c{c}\~oes
e cabos simples. A taxa n\~ao precisa ser alta para um terminal de texto.
Para 115200, o divisor seria 1; para 9600, 12. A escolha n\~ao afeta o
funcionamento do VGA.
\end{nota}

\section{Escrita}

\begin{lstlisting}[caption={Envio de um byte pela UART.},label={lst:uartwrite}]
fn write_byte(byte: u8) {
    unsafe {
        while inb(COM1 + 5) & 0x20 == 0 {}
        outb(COM1, byte);
    }
}
\end{lstlisting}

O la\c{c}o espera o THR esvaziar. Se n\~ao houver UART e o barramento devolver
\code{0xFF}, o bit \code{0x20} j\'a est\'a ligado e o la\c{c}o n\~ao trava;
o byte \'e simplesmente descartado. Se houver uma UART presente mas
travada, o la\c{c}o pode esperar indefinidamente.

\code{write\_str} converte \code{\textbackslash n} em \code{\textbackslash r
\textbackslash n}, porque terminais serial geralmente esperam CRLF:

\begin{lstlisting}[caption={Escrita de string com convers\~ao de nova linha.},label={lst:uartstr}]
pub fn write_str(text: &str) {
    for byte in text.bytes() {
        if byte == b'\n' {
            write_byte(b'\r');
        }
        write_byte(byte);
    }
}
\end{lstlisting}

\section{Leitura}

\begin{lstlisting}[caption={Leitura n\~ao bloqueante de um byte.},label={lst:uartread}]
pub fn read_byte() -> Option<u8> {
    if !PRESENT.load(Ordering::Relaxed) {
        return None;
    }
    unsafe {
        if inb(COM1 + 5) & 0x01 == 0 {
            return None;
        }
        Some(inb(COM1))
    }
}
\end{lstlisting}

A fun\c{c}\~ao \'e n\~ao bloqueante: se n\~ao h\'a byte, retorna \code{None}.
O shell a chama depois de \code{keyboard::read\_char()}, de modo que a
serial funciona como fonte alternativa de entrada.

\section{API}

\begin{table}[H]
\centering
\caption{API de \code{serial.rs}.}
\label{tab:serialapi}
\begin{tabularx}{\linewidth}{@{}lX@{}}
\toprule
\textbf{Fun\c{c}\~ao} & \textbf{Descri\c{c}\~ao} \\
\midrule
\code{init()} & Configura COM1 em 38400 8N1 e detecta presen\c{c}a. \\
\code{read\_byte()} & Retorna um byte recebido ou \code{None}. \\
\code{write\_byte(byte)} & Envia um byte, esperando o THR. Interna. \\
\code{write\_str(text)} & Envia uma string, convertendo \code{\textbackslash n}
em CRLF. \\
\bottomrule
\end{tabularx}
\end{table}

\section{Uso no QEMU e em hardware f\'isico}

No QEMU, a serial pode ser exposta de v\'arias formas:

\begin{lstlisting}[style=shell,caption={Exemplos de exposi\c{c}\~ao da serial no QEMU.},label={lst:qemuserial}]
# Terminal interativo no stdio (usado por `make run`)
qemu-system-x86_64 -cdrom vga-kernel.iso -boot order=d -m 128M -serial stdio

# Socket Unix para testes automatizados
qemu-system-x86_64 -cdrom vga-kernel.iso -boot order=d -m 128M \
  -display none -serial unix:/tmp/vga-serial.sock,server=on,wait=off
\end{lstlisting}

Em um PC f\'isico, COM1 aparece em \code{0x3F8} e pode ser acessada com um
cabo USB-serial e um terminal em 38400 8N1. A BIOS pode desabilitar a porta
ou remape\'a-la; nesse caso o driver n\~ao a encontra.

\section{Limita\c{c}\~oes}

\begin{itemize}
  \item \textbf{Sem IRQs}: a leitura depende de polling.
  \item \textbf{Sem FIFO de entrada}: bytes podem ser perdidos se chegarem
        mais r\'apido do que o shell consulta.
  \item \textbf{Sem controle de fluxo}: RTS/CTS n\~ao s\~ao usados.
  \item \textbf{Sem \emph{echo} autom\'atico}: o pr\'oprio shell ecoa os
        caracteres recebidos.
  \item \textbf{Sem suporte a outras portas}: apenas COM1.
  \item \textbf{Detec\c{c}\~ao heur\'istica}: \code{LSR == 0xFF} pode n\~ao
        ser suficiente em todos os chipsets.
\end{itemize}
